% Generated from locale/tr/content/applied-modal-logic/epistemic-logic/language-epistemic-logic.tex; do not hand-edit.


\olfileid[tr]{aml}{el}{lan}

\olsection{Epistemik Mantığın Dili}

\begin{defn}
$G$ bir özne simgeleri kümesi olsun. Çok özneli epistemik mantığın temel dili
şunları içerir:
\begin{enumerate}
  \tagitem{prvFalse}{!!{falsity} için önerme sabiti~$\lfalse$.}{}
  \tagitem{prvTrue}{!!{truth} için önerme sabiti~$\ltrue$.}{}
  \item !!{denumerable}s bir !!{propositional variable}s kümesi: $\Obj
    p_0$, $\Obj p_1$, $\Obj p_2$, \dots
  \item Önerme bağlaçları: \startycommalist
  \iftag{prvNot}{\ycomma $\lnot$ (değilleme)}{}%
  \iftag{prvAnd}{\ycomma $\land$ (tümel evetleme)}{}%
  \iftag{prvOr}{\ycomma $\lor$ (tikel evetleme)}{}%
  \iftag{prvIf}{\ycomma $\lif$ (!!{conditional})}{}%
  \iftag{prvIff}{\ycomma $\liff$ (!!{biconditional})}{}.
  \item Bilgi işleci $\Knows_a$; burada $a \in G$.
\end{enumerate}
\end{defn}

Sistemimizde yalnızca tek bir öznenin bilgisiyle ilgileniyorsak $G$ kümesine
ve tek tek öznelere yapılan göndermeleri kaldırabiliriz. Bu durumda yalnızca
temel $\Knows$ işlecimiz bulunur.

\begin{defn}
Epistemik dilin \emph{!!^{formula}s{}i} tümevarımlı olarak şöyle tanımlanır:
\begin{enumerate}
\tagitem{prvFalse}{$\lfalse$ atomik bir !!{formula}{}dür.}{}

\tagitem{prvTrue}{$\ltrue$ atomik bir !!{formula}{}dür.}{}

\item Her önerme değişkeni $\Obj p_i$, (atomik) bir !!{formula}{}dür.

\tagitem{prvNot}{$!A$ !!a{formula} ise $\lnot !A$ da
  !!a{formula}{}dür.}{}

\tagitem{prvAnd}{$!A$ ve $!B$ !!{formula}s ise $(!A \land
  !B)$ de !!a{formula}{}dür.}{}

\tagitem{prvOr}{$!A$ ve $!B$ !!{formula}s ise $(!A \lor !B)$
  de !!a{formula}{}dür.}{}

\tagitem{prvIf}{$!A$ ve $!B$ !!{formula}s ise $(!A \lif !B)$
  de !!a{formula}{}dür.}{}

\tagitem{prvIff}{$!A$ ve $!B$ !!{formula}s ise $(!A \liff !B)$
  de !!a{formula}{}dür.}{}

\item $!A$ !!a{formula} ve $a \in G$ ise $\Knows_a !A$ da
  !!a{formula}{}dür.

\tagitem{limitClause}{Başka hiçbir şey !!a{formula} değildir.}{}
\end{enumerate}
\end{defn}

!!a{formula}~$!A$, $\Knows_a$ içermiyorsa ona \emph{modal içermeyen} deriz.

\begin{defn}
  $\Knows$ işleci bireysel bilgiyi simgelemek için kullanılırken, çoğu kez
  ``herkes bilir'' diye okunan $\EKnows$ grup bilgisini simgeler. $G' \subseteq
  G$ olduğunda $\EKnows_{G'} !A$ ifadesini
  \[\bigwedge_{b \in G'} \Knows_b !A\]
  ifadesinin kısaltması olarak tanımlarız.
\end{defn}

Daha da güçlü bir bilgi türünü, yani bir $G$ özne grubu arasındaki
\emph{ortak bilgiyi} de tanımlayabiliriz. Bir bilgi parçasının bir özne grubu
arasında ortak bilgi olması, o gruptaki öznelerin her bileşimi için hepsinin
birbirlerinin birbirlerinin bildiğini \dots sonsuza dek bildikleri anlamına
gelir. Bu, grup bilgisinden önemli ölçüde daha güçlüdür; !!a{formula}{}ün grup
bilgisi olup ortak bilgi olmadığı bağıntısal modeller kurmak kolaydır.
$\CKnows_G !A$ gösterimini ``$G$ grubunda $!A$ ortak bilgidir'' ifadesini
simgelemek için kullanacağız.

